\ifdefined\gkver\else\def\gkver{student}\fi
\documentclass[\gkver]{gaokaozhenti}
\begin{document}

\chapter{1986年全国}

\begin{enumerate}
\item
%% number: 1
%% paperName: 1986年全国高考第 1 题 3 分
%% typeId: 3
%% type: 填空题
%% chapter: 原子和原子核
%% point: 核裂变
%% method: 无
%% score: 3
%% degree: 650
%% duplicateId: 0
%% body:
平衡下列核反应方程：

$\ce{^{235}_{92}U} +$ \tkanswer[1.4cm]{$\ce{^{1}_{0}n}$} $\ce{-> ^{90}_{38}Sr + ^{136}_{54}Xe + 10^{1}_{0}n}$

在核反应堆中，石墨起\tkanswer[2.8cm]{使快中子减速}的作用；镉棒起\tkanswer[3.6cm]{吸收中子控制链式反应速度}的作用。

%% answer: $^{1}_{0}$n；使快中子减速；吸收中子控制链式反应速度
\jdanswer{
$\ce{^{1}_{0}n}$；使快中子减速；吸收中子控制链式反应速度
}
%% memo:
\memoanswer{
【解析】核反应方程满足核电荷数和质量数守恒，由核电荷数守恒有 $38+54+z=92$，质量数守恒有 $90+136+10+y=235$，解得 $z=0$，$y=1$，所以是中子 $\ce{^{1}_{0}n}$。
}

\item
%% number: 2
%% paperName: 1986年全国高考第 2 题 3 分
%% typeId: 3
%% type: 填空题
%% chapter: 光学
%% point: 光的干涉
%% method: 无
%% score: 3
%% degree: 650
%% duplicateId: 0
%% body:
肥皂泡在阳光照射下呈彩色，这是属于光的\tkanswer[1.4cm]{干涉}现象。

低压汞蒸气发光产生的光谱，是属于连续光谱、明线光谱、吸收光谱中的哪一种？答：\tkanswer[2.2cm]{明线光谱}。

无线电波、可见光、伦琴射线、$\gamma$ 射线中的哪一种是原子内层电子受到激发后原子辐射的电磁波？答：\tkanswer[2.2cm]{伦琴射线}。

%% answer: 干涉；明线光谱；伦琴射线
\jdanswer{
干涉；明线光谱；伦琴射线
}
%% memo:
\memoanswer{
本题原卷及材料中未提供详解，待补充。
}

\item
%% number: 3
%% paperName: 1986年全国高考第 3 题 3 分
%% typeId: 3
%% type: 填空题
%% chapter: 曲线运动
%% point: 圆周运动的应用
%% method: 无
%% score: 3
%% degree: 850
%% duplicateId: 0
%% body:
汽车沿半径为 $R$ 的圆跑道行驶，设跑道的路面是水平的，路面作用于车的摩擦力的最大值是车重的 $\frac{1}{10}$，要使车不致冲出圆跑道，车速最大不能超过\tkanswer[2.6cm]{$\sqrt{\frac{gR}{10}}$}。

%% answer: \sqrt{\frac{gR}{10}}
\jdanswer{
$\sqrt{\frac{gR}{10}}$
}
%% memo:
\memoanswer{
【解析】由摩擦力提供向心力得 $\frac{1}{10}mg=m\frac{v^{2}}{R}$，解得 $v=\sqrt{\frac{1}{10}Rg}$。
}

\item
%% number: 4
%% paperName: 1986年全国高考第 4 题 3 分
%% typeId: 3
%% type: 填空题
%% chapter: 光学
%% point: 光电效应
%% method: 无
%% score: 3
%% degree: 750
%% duplicateId: 0
%% body:
某金属用频率为 $\nu_{1}$ 的光照射时产生的光电子的最大初动能，是用频率为 $\nu_{2}$ 的光照射时产生的光电子的最大初动能的 2 倍，则这种金属的逸出功 $W=$\tkanswer[3.2cm]{$h(2\nu_{2}-\nu_{1})$}。

%% answer: h(2ν2−ν1)
\jdanswer{
$h(2\nu_{2}-\nu_{1})$
}
%% memo:
\memoanswer{
【解析】由光电效应方程得 $2E_{\mathrm{k}}=h\nu_{1}-W$，$E_{\mathrm{k}}=h\nu_{2}-W$，联立解得 $W=h(2\nu_{2}-\nu_{1})$。
}

\item
%% number: 5
%% paperName: 1986年全国高考第 5 题 3 分
%% typeId: 3
%% type: 填空题
%% chapter: 光学
%% point: 光的电磁说
%% method: 无
%% score: 3
%% degree: 750
%% duplicateId: 0
%% body:
附图是电台发出的无线电信号的接收电路图（P 为耳机），图中少画了一个元件，请用惯用的符号把这个元件补画在电路图中。图中电容器 $C_{1}$ 起着\tkanswer[2.8cm]{选台（或调频）}的作用。电容器 $C_{2}$ 起着\tkanswer[3.6cm]{通高频阻低频使音频信号通过耳机}的作用。
\onepicture[align=right, w=7cm]{figs/05.png}
%% answer: 补画的元件（二极管）如图所示；选台（或调频）；通高频阻低频使音频信号通过耳机
\jdanswer{
补画的元件（二极管）如图所示；选台（或调频）；通高频阻低频使音频信号通过耳机
}
%% memo:
\memoanswer{
补画的元件（二极管）如图所示；选台（或调频）；通高频阻低频，使音频信号通过耳机。

\onepicture[w=7cm]{figs/05_an.jpg}
}

\item
%% number: 6
%% paperName: 1986年全国高考第 6 题 3 分
%% typeId: 7
%% type: 作图题
%% chapter: 机械波
%% point: 波的图象
%% method: 无
%% score: 3
%% degree: 750
%% duplicateId: 0
%% body:
一列振幅是 $2.0\Ucm$，频率是 $4.0\UHz$ 的简谐横波，以 $32\Ucms$ 的速度沿上图中 $x$ 轴的正方向传播。在某时刻，$x$ 坐标为 $-7.0\Ucm$ 处的介质质点正好经平衡位置且向 $y$ 轴正方向运动。试在图中画出此时刻的波形图（要求至少画出两个波长）。
\onepicture[align=right, w=10cm]{figs/06.png}
%% answer: 如图所示
\jdanswer{
如图所示
}
%% memo:
\memoanswer{
如图所示。

【解析】由题意知 $A=2\Ucm$，$T=\frac{1}{f}=0.25\Us$，由波速公式 $v=\frac{\lambda}{T}$ 得 $\lambda=vT=32\Ucms\times0.25\Us=8\Ucm$。先画出 $x$ 坐标为 $-7.0\Ucm$ 处的介质质点（在平衡位置），然后由该质点向 $y$ 轴正方向运动，判断图象的走向（下坡上行），最后根据波长和振幅画出两个以上波长的波形图。

\onepicture[w=10cm]{figs/06_an.jpg}
}

\item
%% number: 7
%% paperName: 1986年全国高考第 7 题 3 分
%% typeId: 3
%% type: 填空题
%% chapter: 电路
%% point: 故障分析
%% method: 无
%% score: 3
%% degree: 550
%% duplicateId: 0
%% body:
在图示的电路中，灯泡 A 和 B 都是正常发光的。忽然灯泡 B 比原来变暗了些，而灯泡 A 比原来变亮了些。试判断电路中什么地方出现了断路的故障（设只有一处出了故障）。

答：\tkanswer[4.5cm]{$R_{2}$ 所在的支路发生了断路}。
\onepicture[align=right, w=6cm]{\includesvg{figs/07}}
%% answer: R2 所在的支路发生了断路
\jdanswer{
$R_{2}$ 所在的支路发生了断路
}
%% memo:
\memoanswer{
【解析】若 $R_{2}$ 所在的支路发生了断路，则灯泡 A 电路的电阻增大，由电阻串联分压原理知，灯泡 A 电路的电压增大，灯泡 B 电路的电压减小，所以灯泡 B 比原来变暗了些，而灯泡 A 比原来变亮了些。
}

\item
%% number: 8
%% paperName: 1986年全国高考第 8 题 3 分
%% typeId: 3
%% type: 填空题
%% chapter: 电场
%% point: 电场叠加
%% method: 无
%% score: 3
%% degree: 550
%% duplicateId: 0
%% body:
长为 $l$ 的导体棒原来不带电，现将一带电量为 $q$ 的点电荷放在距棒左端 $R$ 处，如图所示。当达到静电平衡后，棒上感应的电荷在棒内中点处产生的场强的大小等于\tkanswer[3.2cm]{$k\frac{q}{(R+\frac{l}{2})^{2}}$}。
\onepicture[align=right, w=7cm]{figs/08.png}
%% answer: k q/(R+l/2)^2
\jdanswer{
$k\frac{q}{(R+\frac{l}{2})^{2}}$
}
%% memo:
\memoanswer{
【解析】带电量为 $q$ 的点电荷在棒内中点处产生的场强为 $E=k\frac{q}{(R+\frac{l}{2})^{2}}$，方向向右。由于在静电平衡时导体中的电场强度为 0，所以棒上感应的电荷在棒内中点处产生的电场场强的大小 $E=k\frac{q}{(R+\frac{l}{2})^{2}}$，方向向左。
}

\item
%% number: 9
%% paperName: 1986年全国高考第 9 题 3 分
%% typeId: 3
%% type: 填空题
%% chapter: 磁场
%% point: 带电粒子在复合场中的运动
%% method: 无
%% score: 3
%% degree: 450
%% duplicateId: 0
%% body:
如下图所示，平行板电容器的极板沿水平方向放置，电子束从电容器左边正中间 a 处沿水平方向入射，电子的初速都是 $v_{0}$，在电场力的作用下，刚好从图中所示的 c 点射出，射出时的速度为 $v$。现若保持电场不变，再加一个匀强磁场，磁场的方向跟电场和电子入射的方向都垂直（图中垂直于纸面向里），使电子刚好由图中 d 点射出，c、d 两点的位置相对于中线 ab 是对称的，则从 d 点射出时每个电子的动能等于\tkanswer[3.2cm]{$mv_{0}^{2}-\frac{1}{2}mv^{2}$}。
\onepicture[align=right, w=7cm]{figs/09.png}
%% answer: mv0^2 − 1/2 mv^2
\jdanswer{
$mv_{0}^{2}-\frac{1}{2}mv^{2}$
}
%% memo:
\memoanswer{
【解析】电子从 a 到 c，电场力做正功，由动能定理得 $W=\frac{1}{2}mv^{2}-\frac{1}{2}mv_{0}^{2}$；电子从 a 到 d，电场力做负功，做功的大小等于从 a 到 c 做的功，又洛伦兹力不做功，由动能定理得 $-W=E_{\mathrm{k}}-\frac{1}{2}mv_{0}^{2}$，解得 $E_{\mathrm{k}}=mv_{0}^{2}-\frac{1}{2}mv^{2}$。
}

\item
%% number: 10
%% paperName: 1986年全国高考第 10 题 4 分
%% typeId: 1
%% type: 单选题
%% chapter: 原子和原子核
%% point: 原子核式结构
%% method: 无
%% score: 4
%% degree: 950
%% duplicateId: 0
%% body:
卢瑟福提出原子的核式结构学说的根据是，在用 $\alpha$ 粒子轰击金箔的实验中发现 $\alpha$ 粒子\xzanswer{B}
\fourchoices[answer=B]
{全部穿过或发生很小的偏转}
{绝大多数穿过，只有少数发生很大偏转，甚至极少数被弹回}
{绝大多数发生很大的偏转，甚至被弹回，只有少数穿过}
{全都发生很大的偏转}
%% answer: B
%% memo:
\memoanswer{
【解析】$\alpha$ 粒子轰击金箔的实验中绝大多数穿过，只有少数发生很大偏转，甚至极少数被弹回，故 B 正确。
}

\item
%% number: 11
%% paperName: 1986年全国高考第 11 题 4 分
%% typeId: 1
%% type: 单选题
%% chapter: 气体性质
%% point: 分子力
%% method: 无
%% score: 4
%% degree: 650
%% duplicateId: 0
%% body:
两个分子甲和乙相距较远（此时它们之间的分子力可忽略），设甲固定不动，乙逐渐向甲靠近直到不能再靠近的整个过程中，\xzanswer{D}
\fourchoices[answer=D]
{分子力总是对乙做正功}
{乙总是克服分子力做功}
{先是乙克服分子力做功，然后分子力对乙做正功}
{先是分子力对乙做正功，然后乙克服分子力做功}
%% answer: D
%% memo:
\memoanswer{
【解析】因为存在一个距离 $r_{0}$，在分子间距离较远时（$r>r_{0}$），分子力表现为引力，所以乙逐渐向甲靠近过程中分子力对乙做正功；当分子间距离较近时（$r<r_{0}$），分子力表现为斥力，所以乙逐渐向甲靠近过程中分子力对乙做负功，即乙克服分子力做功，故 D 正确。
}

\item
%% number: 12
%% paperName: 1986年全国高考第 12 题 4 分
%% typeId: 1
%% type: 单选题
%% chapter: 力物体的平衡
%% point: 物体系平衡
%% method: 无
%% score: 4
%% degree: 550
%% duplicateId: 0
%% body:
如图所示，一个箱子放在水平地面上，箱内有一固定的竖直杆，在杆上套着一个环。箱和杆的质量为 $M$，环的质量为 $m$。已知环沿着杆加速下滑，环与杆的摩擦力的大小为 $f$，则此时箱对地面的压力\xzanswer{C}
\onepicture[w=5cm]{figs/12.png}
\fivechoices[answer=C]
{等于 $Mg$}
{等于（$M+m$）$g$}
{等于 $Mg+f$}
{等于（$M+m$）$g-f$}
{无法确定}
%% answer: C
%% memo:
\memoanswer{
【解析】对环，受重力和摩擦力的作用，摩擦力大小为 $f$，方向向上；对杆，受摩擦力大小也为 $f$，方向向下；对箱和杆，受重力 $Mg$ 和支持力 $N$ 的作用，由力的平衡得 $Mg+f=N$，故 C 正确。
}

\item
%% number: 13
%% paperName: 1986年全国高考第 13 题 4 分
%% typeId: 1
%% type: 单选题
%% chapter: 运动和力
%% point: 牛顿定律的应用
%% method: 无
%% score: 4
%% degree: 550
%% duplicateId: 0
%% body:
在有空气阻力的情况下，以初速 $v_{1}$ 竖直上抛一物体，经过时间 $t_{1}$ 到达最高点。又经过时间 $t_{2}$，物体由最高点落回到抛出点，这时物体的速度为 $v_{2}$。则\xzanswer{E}
\fivechoices[answer=E]
{$v_{2}=v_{1}$，$t_{2}=t_{1}$}
{$v_{2}>v_{1}$，$t_{2}>t_{1}$}
{$v_{2}<v_{1}$，$t_{2}<t_{1}$}
{$v_{2}>v_{1}$，$t_{2}<t_{1}$}
{$v_{2}<v_{1}$，$t_{2}>t_{1}$}
%% answer: E
%% memo:
\memoanswer{
【解析】由于物体克服空气阻力做功，所以速度 $v_{2}<v_{1}$。设空气阻力为 $f$，在上升阶段，加速度为 $a_{1}=g+\frac{f}{m}$，在下落阶段，加速度为 $a_{2}=g-\frac{f}{m}$，位移为 $h=\frac{1}{2}at^{2}$，所以 $t=\sqrt{\frac{2h}{a}}$。因为 $a_{1}>a_{2}$，所以 $t_{2}>t_{1}$，故 E 正确。
}

\item
%% number: 14
%% paperName: 1986年全国高考第 14 题 4 分
%% typeId: 2
%% type: 多选题
%% chapter: 电场
%% point: 电场线
%% method: 无
%% score: 4
%% degree: 750
%% duplicateId: 0
%% body:
指出下图所示的哪些情况中，a、b 两点的电势相等，a、b 两点的电场强度矢量也相等。\xzanswer{BD}
\fourpicture[numA=A, numB=B, numC=C, numD=D, w=2.6cm]
{figs/14a.png}
{figs/14b.png}
{figs/14c.png}
{figs/14d.png}
\fourchoices[answer=BD]
{平行板电容器带电时，极板间除边缘附近外的任意两点 a、b}
{静电场中达到静电平衡时的导体内部任意两点 a、b}
{离点电荷等距的任意两点 a、b}
{两个等量异号的点电荷，在其连线的中垂线上，与连线中点 O 等距的两点 a、b}
%% answer: BD
%% memo:
\memoanswer{
【解析】对 A 图，a、b 两点的电场强度矢量相等，电势不等；对 B 图，a、b 两点的电场强度矢量相等，电势也相等；对 C 图，a、b 两点的电场强度矢量不等（大小相等，方向不同），电势相等；对 D 图，a、b 两点的电场强度矢量相等，电势相等（为 0），故 BD 正确。
}

\item
%% number: 15
%% paperName: 1986年全国高考第 15 题 4 分
%% typeId: 2
%% type: 多选题
%% chapter: 电场
%% point: 带电粒子的圆周运动
%% method: 无
%% score: 4
%% degree: 450
%% duplicateId: 0
%% body:
如图所示是氢原子中电子绕核做快速的圆周运动（设为逆时针）的示意图。电子绕核运动，可等效为环形电流。设此环形电流在通过圆心并垂直于圆面的轴线上某一点 P 处产生的磁感应强度的大小为 $B_{1}$。现在沿垂直于圆轨道平面的方向加一磁感应强度为 $B_{0}$ 的外磁场，这时设电子的轨道半径没变，而它的速度发生了变化。若用 $B_{2}$ 表示此时环形电流在 P 点产生的磁感应强度的大小，则当 $B_{0}$ 的方向\xzanswer{BC}
\onepicture[w=3.4cm]{figs/15.png}
\fivechoices[answer=BC]
{垂直于纸面向里时，$B_{2}>B_{1}$}
{垂直于纸面向里时，$B_{2}<B_{1}$}
{垂直于纸面向外时，$B_{2}>B_{1}$}
{垂直于纸面向外时，$B_{2}<B_{1}$}
{不论是垂直于纸面向里还是向外，$B_{1}$ 总是等于 $B_{2}$}
%% answer: BC
%% memo:
\memoanswer{
【解析】①当外加磁场 $B_{0}$ 的方向垂直于纸面向里时，由左手定则，洛伦兹力方向指向圆心外侧，库仑力与洛伦兹力之差提供向心力，有 $F=m\frac{v^{2}}{r}$。因 $r$ 不变，所以速度 $v$ 随向心力 $F$ 的减小而减小，故等效环形电流减小，环形电流产生的磁场的磁感应强度的大小也减小，即 $B_{2}<B_{1}$，故 B 正确；②当外加磁场 $B_{0}$ 的方向垂直于纸面向外时，由左手定则，洛伦兹力方向指向圆心，库仑力与洛伦兹力之和提供向心力，有 $F=m\frac{v^{2}}{r}$。因 $r$ 不变，所以速度 $v$ 随向心力 $F$ 的增大而增大，故等效环形电流增大，环形电流产生的磁场的磁感应强度的大小也增大，即 $B_{2}>B_{1}$，故 C 正确。
}

\item
%% number: 16
%% paperName: 1986年全国高考第 16 题 4 分
%% typeId: 1
%% type: 单选题
%% chapter: 直线运动
%% point: 相遇问题
%% method: 图象法
%% score: 4
%% degree: 550
%% duplicateId: 0
%% body:
汽车甲沿着平直的公路以速度 $v_{0}$ 做匀速直线运动。当它路过某处的同时，该处有一辆汽车乙开始做初速为零的匀加速运动去追赶甲车。根据上述的已知条件\xzanswer{A}
\fourchoices[answer=A]
{可求出乙车追上甲车时乙车的速度}
{可求出乙车追上甲车时乙车所走的路程}
{可求出乙车从开始起动到追上甲车时所用的时间}
{不能求出上述三者中任何一个}
%% answer: A
%% memo:
\memoanswer{
公式法：$v_{0}t=\frac{1}{2}at^{2}$，因 $a$ 未知，不能求出 $t$，也不能求出 $s$；由 $v=at=2v_{0}$ 可求出乙车追上甲车时乙车的速度，故 A 正确；

图像法：从图像可以看出，加速度不同，时间不同，位移也不同，但末速度相同，都是 $v=2v_{0}$。

\onepicture[align=right, w=5.5cm]{figs/16_an.png}
}

\item
%% number: 17
%% paperName: 1986年全国高考第 17 题 7 分
%% typeId: 4
%% type: 实验题
%% chapter: 电路
%% point: 伏安法测电阻
%% method: 无
%% score: 7
%% degree: 750
%% duplicateId: 0
%% body:
有一电阻 $R_{x}$，其阻值大约在 $40\sim50$ 欧姆之间。需要进一步测定其阻值，手边现有下列器材：

\begin{center}
\begin{tabular}{|c|c|c|}
\hline
器材 & 代号 & 规格 \\
\hline
电池组 & $E$ & 电动势 $9\UV$，内阻约 $0.5\UO$ \\
\hline
伏特表 & V & 量程 $0\sim10\UV$，内阻 $20\UkO$ \\
\hline
毫安表 & $A_{1}$ & 量程 $0\sim50\UmA$，内阻约 $20\UO$ \\
\hline
毫安表 & $A_{2}$ & 量程 $0\sim300\UmA$，内阻约 $4\UO$ \\
\hline
滑动变阻器 & $R_{1}$ & 阻值范围 $0\sim100\UO$，额定电流 $1\UA$ \\
\hline
滑动变阻器 & $R_{2}$ & 阻值范围 $0\sim1700\UO$，额定电流 $0.3\UA$ \\
\hline
电键 & K & \\
\hline
几根连接用导线 & & \\
\hline
\end{tabular}
\end{center}

有两种可供选用的电路如图 \ref{1986全国17a} 和图 \ref{1986全国17b} 所示。

实验要求多测几组电流、电压值，画出电流$-$电压关系图。

为了实验能正常进行并减小测量误差，而且要求滑动变阻器便于调节，在实验中应选图\tkanswer[1cm]{1}所示的电路，应选代号是\tkanswer[1.4cm]{$A_{2}$}的毫安表和代号是\tkanswer[1.4cm]{$R_{1}$}的滑动变阻器。

\twopicture[numA=1, numB=2, labelA=1986全国17a, labelB=1986全国17b, align=right, w=6.5cm]
{figs/17a.png}
{figs/17b.png}
%% answer: 1；A2；R1
\jdanswer{
1；$A_{2}$；$R_{1}$
}
%% memo:
\memoanswer{
题中图 1 是电流表外接法，图 2 是电流表内接法。因为电压表的内阻比待测电阻大得多，$\frac{20000}{50}=400$ 倍，而待测电阻比电流表内阻大得不多，$\frac{40}{4}=10$ 倍，所以用电流表外接法，选图 1；电路中的电流约为 $I=\frac{E}{R}=\frac{9}{50}\UA=0.2\UA=200\UmA$，$A_{1}$ 的量程不能满足要求，电流表选 $A_{2}$，因电流表外接，内阻大小不是主要因素；滑动变阻器为便于调节，选 $R_{1}$。
}

\item
%% number: 18
%% paperName: 1986年全国高考第 18 题 3 分
%% typeId: 4
%% type: 实验题
%% chapter: 直线运动
%% point: 游标卡尺和螺旋测微仪
%% method: 无
%% score: 3
%% degree: 850
%% duplicateId: 0
%% body:
用螺旋测微器测量一矩形小零件的长和宽时，螺旋测微器上的示数如图 \ref{1986全国18a} 和图 \ref{1986全国18b} 所示。图 \ref{1986全国18a} 的读数是\tkanswer[1.6cm]{8.474}$\Umm$。图 \ref{1986全国18b} 的读数是\tkanswer[1.6cm]{6.576}$\Umm$。

\twopicture[numA=1, numB=2, labelA=1986全国18a, labelB=1986全国18b, align=right, w=5cm]
{figs/18a.png}
{figs/18b.png}
%% answer: 8.474；6.576（允许最后一位差±1）
\jdanswer{
$8.474\Umm$；$6.576\Umm$（允许最后一位差 $\pm1$）
}
%% memo:
\memoanswer{
图 1 中，固定刻度读 $8\Umm$，可动刻度读 $47.4\times0.01\Umm=0.474\Umm$，故图 1 的读数为 $8.474\Umm$；图 2 中，固定刻度读 $6.5\Umm$，可动刻度读 $7.6\times0.01\Umm=0.076\Umm$，故图 2 的读数为 $6.576\Umm$。最后一位允许差 $\pm1$。
}

\item
%% number: 19
%% paperName: 1986年全国高考第 19 题 4 分
%% typeId: 4
%% type: 实验题
%% chapter: 气体性质
%% point: 验证玻意耳定律
%% method: 无
%% score: 4
%% degree: 650
%% duplicateId: 0
%% body:
一个学生用带有刻度的注射器做验证玻意耳$-$马略特定律的实验，他做实验时的主要步骤如下：

\begin{steps}
	\item
	用刻度尺测出注射器全部刻度之长，用这个长度去除它的容积得出活塞的横截面积 $S$。
	\item
	用天平称出活塞的质量 $M$。
	\item
	把适量的润滑油均匀地抹在注射器的活塞上，把活塞插进注射器内一部分，然后将注射器的小孔用橡皮帽堵住，记下这时空气柱的体积 $V$。
	\item
	用烧瓶夹把注射器竖直固定在铁架上，利用砝码重量向下压活塞，使空气柱体积减小。改变砝码个数，再做两次，记下每次砝码的质量 $m$ 和相应的空气柱的体积 $V$。
	\item
	把记录的数据填入表格里，根据算式 $p=\frac{(M+m)g}{S}$ 计算出每次的压强值。
	\item
	求出每次压强 $p$ 跟相应的体积 $V$ 的乘积，看看它们是否相等。
\end{steps}

根据你做这个实验的经验，这个学生的实验步骤中有些什么重要错误或疏漏。答：\tkanswer[6cm]{步骤 5 中所用的公式是错误的，应改成 $p=\frac{(M+m)g}{S}+p_{0}$，$p_{0}$ 为大气压；在进行步骤 5 之前要读气压计所指示的大气压强 $p_{0}$ 的值}。

%% answer: 步骤 5 中所用的公式是错误的，应改成 p=(M+m)g/S+p0，p0 为大气压；在进行步骤 5 之前要读气压计所指示的大气压强 p0 的值
\jdanswer{
步骤 5 中所用的公式是错误的，应改成 $p=\frac{(M+m)g}{S}+p_{0}$，$p_{0}$ 为大气压。在进行步骤 5 之前要读气压计所指示的大气压强 $p_{0}$ 的值。
}
%% memo:
\memoanswer{
步骤 5 中所用的公式是错误的，应改成 $p=\frac{(M+m)g}{S}+p_{0}$，$p_{0}$ 为大气压。在进行步骤 5 之前要读气压计所指示的大气压强 $p_{0}$ 的值。
}

\item
%% number: 20
%% paperName: 1986年全国高考第 20 题 8 分
%% typeId: 6
%% type: 计算题
%% chapter: 电磁感应
%% point: 电磁感应受力分析
%% method: 无
%% score: 8
%% degree: 850
%% duplicateId: 0
%% body:
图中 abcd 是一个固定的 U 形金属框架，ab 和 cd 边都很长，bc 边长为 $l$，框架的电阻可不计，ef 是放置在框架上与 bc 平行的导体杆，它可在框架上自由滑动（摩擦可忽略），它的电阻为 $R$。现沿垂直于框架平面的方向加一恒定的匀强磁场，磁感应强度为 $B$，方向垂直于纸面向里。已知当以恒力 $F$ 向右拉导体杆 ef 时，导体杆最后匀速滑动。求匀速滑动时的速度。
\onepicture[align=right, w=6cm]{figs/20.png}
%% answer: v=FR/(B^2 l^2)
\jdanswer{
$v=\frac{FR}{B^{2}l^{2}}$
}
%% memo:
\memoanswer{
当导体杆向右滑动时，通过回路 ebcf 的磁通量将发生变化，从而在回路中产生感生电动势 $\varepsilon$ 和感生电流 $I$。设导体杆做匀速运动时的速度为 $v$，根据法拉第电磁感应定律和欧姆定律可知

$\varepsilon=Blv$　（a）

$I=\frac{E}{R}$　（b）

根据安培定律可知，磁场对导体杆的作用力为

$f=BIl$　（c）

方向向左，即与外力 $F$ 的方向相反。当导体杆做匀速运动时，有

$f=F$　（d）

由以上可解得 $v=\frac{FR}{B^{2}l^{2}}$。

评分标准：本题 8 分。列出（a）式的，给 2 分；列出（b）式的，给 2 分；列出（c）式的，给 2 分；列出（d）式的，给 1 分。最后结果正确的，再给 1 分。
}

\item
%% number: 21
%% paperName: 1986年全国高考第 21 题 8 分
%% typeId: 6
%% type: 计算题
%% chapter: 运动和力
%% point: 动量和能量
%% method: 无
%% score: 8
%% degree: 550
%% duplicateId: 0
%% body:
甲、乙两个小孩各乘一辆冰车在水平冰面上游戏。甲和他的冰车的质量共为 $M=30\Ukg$，乙和他的冰车的质量也是 $30\Ukg$。游戏时，甲推着一个质量为 $m=15\Ukg$ 的箱子，和他一起以大小为 $v_{0}=2.0\Ums$ 的速度滑行，乙以同样大小的速度迎面滑来。为了避免相撞，甲突然将箱子沿冰面推给乙，箱子滑到乙处时乙迅速把它抓住。若不计冰面的摩擦力，求：

\begin{enumerate}
	\item
	甲至少要以多大的速度（相对于地面）将箱子推出，才能避免与乙相撞。
	\item
	甲在推出时对箱子做了多少功。
\end{enumerate}
\onepicture[align=right, w=9cm]{figs/21.png}
%% answer: （1）5.2 m/s；（2）1.7×10^2 J
\jdanswer{
\begin{enumerate}
	\item
	$v=5.2\Ums$。

	\item
	$W=1.7\times10^{2}\UJ$。
\end{enumerate}
}
%% memo:
\memoanswer{
（1）在推出和抓住的过程中，小孩、冰车和箱子的总动量守恒。要想刚能避免相碰，要求抓住后甲和乙的速度正好相等。由此就可求得推出时的最小速度。

设箱子推出后其速度为 $v$，甲孩的速度为 $v_{1}$，根据动量守恒可得

$mv+Mv_{1}=(m+M)v_{0}$　（a）

设乙孩抓住箱子后其速度为 $v_{2}$，根据动量守恒可得

$(m+M)v_{2}=mv-Mv_{0}$　（b）

刚好不相碰的条件要求

$v_{1}=v_{2}$　（c）

由（a）、（b）、（c）三式可解得

$v=\left(\frac{m^{2}+2mM+2M^{2}}{m^{2}+2mM}\right)v_{0}$

代入数值可得 $v=5.2\Ums$。

（2）设推出时甲对箱子做功为 $W$，根据功能关系可知

$W=\frac{1}{2}mv^{2}-\frac{1}{2}mv_{0}^{2}$　（d）

代入数值可得

$W=1.7\times10^{2}\UJ$。

评分标准：本题 8 分。（1）占 6 分；（2）占 2 分。

（1）中，列出（a）式的，给 1 分；列出（b）式的，给 1 分；列出（c）式的，给 3 分。答数正确的，再给 1 分。

（2）中，列出（d）式的，给 1 分；答数正确的，再给 1 分。
}

\item
%% number: 22
%% paperName: 1986年全国高考第 22 题 9 分
%% typeId: 6
%% type: 计算题
%% chapter: 力物体的平衡
%% point: 一般平衡条件
%% method: 无
%% score: 9
%% degree: 450
%% duplicateId: 0
%% body:
一个质量为 $m=50\Ukg$ 的均匀圆柱体，放在台阶的旁边，台阶的高度 $h$ 是柱体半径 $r$ 的一半，如图所示（图为其横截面），柱体与台阶接触处（图中 P 点所示）是粗糙的。现要在图中柱体的最上方 A 处施一最小的力，使柱体刚能开始以 P 为轴向台阶上滚，求：

\begin{enumerate}
	\item
	所加的力的大小；
	\item
	台阶对柱体的作用力的大小。
\end{enumerate}
\onepicture[align=right, w=5cm]{figs/22.png}
%% answer: （1）F=2.5×10^2 N；（2）f=4.3×10^2 N
\jdanswer{
\begin{enumerate}
	\item
	$F=2.5\times10^{2}\UN$

	\item
	$f=4.3\times10^{2}\UN$
\end{enumerate}
}
%% memo:
\memoanswer{
（1）要在 A 处施一最小的力，则力的方向应与 AP 垂直，这样力臂最大。因为 $r=2h$，由几何关系可推知 $\angle PAO=30^{\circ}$，$\angle POB=60^{\circ}$。要使柱体刚能绕 P 轴上滚，即意味着此时地面对柱体的支持力

$N=0$　（a）

这时，拉力 $F$ 和重力 $mg$ 对 P 轴的力矩平衡，由此可得

$mgr\sin60^{\circ}=F\cdot2r\cos30^{\circ}$　（b）

\onepicture[w=6cm]{figs/22_an.jpg}

所以 $F=2.5\times10^{2}\UN$。

（2）设台阶对柱体的作用力为 $f$，因为刚能开始运动时，$f$ 与重力 $mg$ 及拉力 $F$ 三力平衡，所以必为共点力。由此可知力 $f$ 的方向是沿 PA 方向。即力 $f$ 的方向与 $F$ 的方向垂直，所以 $f$ 的大小必等于重力在 AP 方向上的分力。即

$f=mg\cos30^{\circ}$　（c）

$f=4.3\times10^{2}\UN$。

评分标准：本题 9 分。（1）占 6 分；（2）占 3 分。

（1）中，知道 F 与 AP 垂直的，给 3 分；列出（a）式的，给 1 分；列出（b）式的，给 1 分；答数正确的，再给 1 分。

（2）中，列出（c）式的，给 2 分；答数正确的，再给 1 分。
}

\item
%% number: 23
%% paperName: 1986年全国高考第 23 题 6 分
%% typeId: 6
%% type: 计算题
%% chapter: 电路
%% point: 闭合电路欧姆定律
%% method: 无
%% score: 6
%% degree: 450
%% duplicateId: 0
%% body:
有两只伏特表 A 和 B，量程已知，内阻不知等于多少。另有一干电池，它的内阻不能忽略，但不知等于多少。只用这两只伏特表、电键和一些连接用导线，能通过测量计算出这个电池的电动势（已知电动势不超出伏特表的量程，干电池不许拆开）。

\begin{enumerate}
	\item
	画出你测量时所用的电路图。
	\item
	以测得的量做为已知量，导出计算电动势的式子。
\end{enumerate}
%% answer: （1）测量时用的电路图如图所示；（2）\varepsilon=\frac{U_A U_B}{U_A-U_A'}
\jdanswer{
\begin{enumerate}
	\item
	测量时用的电路图如图所示。

	\item
	$\varepsilon=\frac{U_{\mathrm{A}}U_{\mathrm{B}}}{U_{\mathrm{A}}-U_{\mathrm{A}}'}$
\end{enumerate}
}
%% memo:
\memoanswer{
（1）测量时用的电路图如图所示。

\onepicture[w=10cm]{figs/23_an.jpg}

（2）将任一只伏特表，如表 A，与电源连通，记下伏特表所示的电压值 $U_{A}$。再将两只伏特表与电源串联，记下伏特表所示的电压值 $U_{A}'$ 和 $U_{B}$。设表 A 的内阻为 $R_{A}$，表 B 的内阻为 $R_{B}$，电源内阻为 $r$，电源电动势为 $\varepsilon$，只将表 A 与电源连通时电流强度为 $I_{1}$，两表与电源串联时电流强度为 $I_{2}$，根据欧姆定律可得

$\varepsilon=I_{1}(R_{A}+r)$　（a）

$\varepsilon=I_{2}(R_{A}+R_{B}+r)$　（b）

因为

$U_{A}=I_{1}R_{A},\quad U_{A}'=I_{2}R_{A},\quad U_{B}=I_{2}R_{B},$　（c）

所以由以上诸式可得

$\varepsilon=\frac{U_{\mathrm{A}}U_{\mathrm{B}}}{U_{\mathrm{A}}-U_{\mathrm{A}}'}$　（d）

评分标准：本题 6 分。（1）占 2 分；（2）占 4 分。

（1）中，只画出一个电路图的，给 0 分。

（2）中，列出（a）式和（b）式的，给 1 分，只列出其中一个的，给 0 分；列出（c）式的，给 1 分；得出（d）式的，再给 2 分。
}

\item
%% number: 24
%% paperName: 1986年全国高考第 24 题 10 分
%% typeId: 6
%% type: 计算题
%% chapter: 电路
%% point: 电容
%% method: 无
%% score: 10
%% degree: 550
%% duplicateId: 0
%% body:
有一个平行板电容器，当两极板间为空气时，其电容为 $C_{0}=40\UpF$，把它连接到一个电动势为 $\varepsilon=500\UV$ 的电源上。现将一块厚度等于极板间距离的石蜡块塞进两极板间，使它充满极板间空间的一半，如图。已知石蜡的介电常数 $\varepsilon=2$。求：

\begin{enumerate}
	\item
	塞入石蜡块后，电容器的电容 $C$；
	\item
	在石蜡块塞入过程中，电源所提供的电能。
\end{enumerate}
\onepicture[align=right, w=4cm]{figs/24.png}
%% answer: （1）C=60 pF；（2）W=5.0×10^-6 J
\jdanswer{
\begin{enumerate}
	\item
	$C=60\UpF$

	\item
	$W=5.0\times10^{-6}\UJ$
\end{enumerate}
}
%% memo:
\memoanswer{
（1）石蜡塞入后，电容器可视为两个平行板电容器并联。每个的极板间距离都是 $d$，每个的极板面积都是原电容器的极板面积 $S$ 的一半。因此，极板间是空气的那一半的电容

$C_{1}=\frac{\frac{S}{2}}{4\pi kd}=\frac{C_{0}}{2}$　（a）

极板间是石蜡的那一半的电容

$C_{2}=\frac{\frac{\xi S}{2}}{4\pi kd}=\frac{\xi C_{0}}{2},$　（b）

所以电容器的总电容

$C=C_{1}+C_{2}$　（c）

由（a）、（b）、（c）三式，代入数值可解得

$C=60\UpF$。

（2）电源提供的能量 $W$ 应等于电动势和通过电源的电量的乘积。设插入石蜡前极板上的电量为 $Q_{0}$，插入后为 $Q$，则

$W=\varepsilon(Q-Q_{0})$。　（d）

因为

$Q_{0}=C_{0}\varepsilon,\ Q=C\varepsilon,$　（e）

所以得

$W=\varepsilon^{2}(C-C_{0})$。

代入数值得

$W=5.0\times10^{-6}\UJ$。

评分标准：本题 10 分。（1）占 5 分；（2）占 5 分。

（1）中，列出（a）式的，给 1 分；列出（b）式的，给 1 分；列出（c）式的，给 2 分；答数正确的，再给 1 分。

（2）中，列出（d）式的，给 3 分；列出（e）式的，给 1 分；答数正确的，再给 1 分。
}

\end{enumerate}

\end{document}
